% Source: https://artofproblemsolving.com/community/c3h2378595_question_on_divisibility
% Stable identity: aops-topic-2378595-q01
Since $111=3\cdot37$ with $3$ and $37$ prime, we need $3\mid16^n-1$ and $37\mid16^n-1$.

\textbf{Modulo 3.} $16\equiv1$, so $16^n\equiv1\pmod3$ for every $n$.

\textbf{Modulo 37.} Compute powers of $16$:
\[16^2=256=6\cdot37+34\equiv-3,\qquad16^3\equiv16\cdot(-3)=-48\equiv-11,\qquad16^9\equiv(-11)^3=-1331\equiv1,\]
since $1331=36\cdot37-1$. The order of $16$ modulo $37$ (the least $d\ge1$ with $16^d\equiv1$) divides $9$, and it is not $1$ or $3$ (as $16\not\equiv1$ and $16^3\equiv-11\not\equiv1$), so it is $9$. Hence $16^n\equiv1\pmod{37}$ exactly when $9\mid n$: writing $n=9q+r$ with $0\le r<9$ gives $16^n\equiv16^r$, which is $1$ only for $r=0$.

So $111\mid16^n-1$ exactly when $9\mid n$. The two-digit multiples of $9$ are $18,27,\dots,99$: there are $\boxed{10}$ of them.
